\clearpage\section{Restricted native progress and complete prose proofs}\label{app:native}
This appendix reproduces the mathematical parts of the 8 September 2026
native progress note. The polynomial obstruction, finite reconstruction,
rectangularity and topological arguments are prose proofs. The intrinsic
commutative-algebra obstruction is the stated classical Stacks specialization,
with its original source record and finiteness argument. None of these prose
arguments is promoted to a new Lean run or to a completion of unrestricted
Native 13.2 or 13.3.
\hypertarget{native:a-complete-polynomial-obstruction-theorem}{%
\subsection{2. A complete polynomial obstruction
theorem}\label{native:a-complete-polynomial-obstruction-theorem}}

\textbf{Scope.} Fix a prime \(p\), an integer \(L\ge1\), finite natural
numbers \(d,r\), and polynomials
\(F_1,\ldots,F_r\in\mathbb Z[X_1,\ldots,X_d]\). Let
\(\bar x\in(\mathbb Z/p^L\mathbb Z)^d\) satisfy every equation
\(F_i(\bar x)=0\). Choose any integer representative
\(x\in\mathbb Z^d\). Set

\[e_i=F_i(x)/p^L\pmod p,\qquad J_{ij}=\frac{\partial F_i}{\partial X_j}(x)\pmod p.\]

The divisions are exact integer divisions. Regard \(J\) as a linear map
\(\mathbb F_p^d\to\mathbb F_p^r\). Define

\[O_{F,\bar x}=\mathbb F_p^r/\operatorname{im}J,\qquad o_{F,\bar x}=[e].\]

\textbf{Theorem.} For this fixed polynomial presentation, the group and
distinguished class are independent of the chosen integer
representative. The class vanishes if and only if \(\bar x\) lifts to a
solution modulo \(p^{L+1}\). The obstruction group has
\(p^{r-\operatorname{rank}J}\) elements. When lifts exist, their set is
a torsor for \(\ker J\) and contains exactly
\(p^{d-\operatorname{rank}J}\) elements.

\textbf{Proof.} For every integer correction vector \(t\), expansion of
each monomial gives

\[F(x+p^Lt)\equiv F(x)+p^LJ(x)t\pmod {p^{L+1}}.\tag{1}\]

Indeed, each term involving at least two correction factors is divisible
by \(p^{2L}\), and \(2L\ge L+1\). Constant and linear correction terms
give the two displayed summands. Adding the monomial expansions proves
(1) for every polynomial, with no denominator or analytic convergence
issue.

Every element reducing to \(\bar x\) is uniquely expressible modulo
\(p^{L+1}\) as \(x+p^Lt\) with \(t\in\mathbb F_p^d\). By (1), it solves
the equations exactly when

\[e+Jt=0\quad\text{in }\mathbb F_p^r.\tag{2}\]

Equation (2) is soluble exactly when \(e\in\operatorname{im}J\): the
sign makes no difference because the image is a subspace. This is
equivalent to \([e]=0\).

Another representative has the form \(x'=x+p^Ls\). Since \(L\ge1\), its
Jacobian is the same modulo \(p\). Dividing (1) by \(p^L\) shows that
its residual digit is \(e'=e+Js\), so it gives exactly the same quotient
class. If \(t_0\) solves (2), all its solutions are \(t_0+\ker J\), and
no others. The finite cardinalities follow from rank-nullity. This also
covers no variables, no equations, zero Jacobian and singular solutions.
\(\square\)

The proof is a universal mathematical argument. The accompanying Python
checks are independently bounded evidence about an implementation, not a
substitute for that argument. The standalone Lean specialization does
not formalize multivariable polynomial expansion (1), unless its own
checked source explicitly says so.

\hypertarget{native:necessary-boundaries-and-exact-examples}{%
\subsubsection{Necessary boundaries and exact
examples}\label{native:necessary-boundaries-and-exact-examples}}

\begin{itemize}
\tightlist
\item
  \textbf{A failed singular lift:} \(p=3\), \(L=1\), \(F(X)=X^2-3\),
  \(x=0\). Here \(e=-1\ne0\) in \(\mathbb F_3\) and \(J=0\). There are
  no lifts of zero modulo 3 to a root modulo 9.
\item
  \textbf{Representative independence is a quotient assertion:} for
  \(F(X)=X^2-1\), \(p=3\), \(L=1\), representatives \(1\) and \(4\) have
  residual digits \(0\) and \(2\). Their common Jacobian is \(2\), with
  full image, and both give zero obstruction. Their corrected
  representatives \(1\) and \(10\) are the same root modulo 9.
\item
  \textbf{A successful step need not give indefinite continuation:} for
  \(F(X)=X^2-5\), \(p=2\), the root \(1\) modulo 2 has zero obstruction
  and lifts to \(1,3\) modulo 4. Both have nonzero obstruction at the
  next step. There is no root modulo 8 and no compatible infinite lift.
\item
  \textbf{The raw Jacobian cokernel need not be presentation
  independent:} the presentations \(F=(X)\) and \(G=(X,0)\) define the
  same solution system. Their raw obstruction groups at zero are
  respectively zero and \(\mathbb F_p\). Both actual distinguished
  obstructions are zero. Thus one cannot claim a canonical isomorphism
  of these two raw groups merely because the equations define the same
  seed.
\end{itemize}

\hypertarget{native:an-intrinsic-finite-obstruction-under-an-algebraic-restriction}{%
\subsection{3. An intrinsic finite obstruction under an algebraic
restriction}\label{native:an-intrinsic-finite-obstruction-under-an-algebraic-restriction}}

Let \(B\) be a finitely presented commutative \(\mathbb Z\)-algebra,
\(p\) prime and \(L\ge1\). For a ring map
\(f:B\to\mathbb Z/p^L\mathbb Z\), use its reduction to make
\(\mathbb F_p\) a \(B\)-module. There is a canonical class

\[o(f)\in\operatorname{Ext}^1_B(\mathrm{NL}_{B/\mathbb Z},\mathbb F_p)\]

whose vanishing is equivalent to a lift
\(B\to\mathbb Z/p^{L+1}\mathbb Z\). This follows from the ring-map
lifting theorem in
\href{https://stacks.math.columbia.edu/tag/08S5}{Stacks Project, Lemma
91.2.1, tag 08S5}. Specialize the source extension to
\(0\to0\to B\to B\) and the target extension to the kernel
\(p^L/p^{L+1}\cong\mathbb F_p\). The base maps are identities of
\(\mathbb Z\). The kernel is square zero because \(L\ge1\).

The comparison in
\href{https://stacks.math.columbia.edu/tag/00S1}{Stacks Project, Lemma
10.134.2, tag 00S1} identifies the complex from a polynomial
presentation with the canonical naive cotangent complex up to homotopy.
Thus this construction is intrinsic, up to the stated canonical
comparison.

The group is finite. Choose a finite presentation
\(P=\mathbb Z[X_1,\ldots,X_d]\twoheadrightarrow B\) with finitely
generated kernel \(I\). The presentation description in
\href{https://stacks.math.columbia.edu/tag/0GPU}{Stacks Project, Remark
91.2.4, tag 0GPU} identifies the group as

\[\operatorname{coker}\left(\operatorname{Hom}_B(\Omega_{P/\mathbb Z}\otimes_P B,\mathbb F_p)\longrightarrow\operatorname{Hom}_B(I/I^2,\mathbb F_p)\right).\]

If \(I\) has \(r\) chosen generators, evaluation on their classes
injects the latter Hom module into \(\mathbb F_p^r\). Hence its quotient
is finite. The relation module \(I/I^2\) retains dependencies between
equations, including redundant zero equations, which the raw target
\(\mathbb F_p^r\) does not.

\textbf{Status:} externally proved after the exact stated
specialization, with the finiteness argument above. This theorem is not
independently formalized in the accompanying Lean file. Extending it to
all operator/recurrence/module systems admitted by manuscript Definition
1.1 requires a specified encoding and a proof that its notion of lift is
preserved. That is the missing implication for applying it to native
Conjecture 13.2 without restriction.

\textbf{Source record.} The Stacks Project authors, live primary text
inspected 2026-09-08, tags 08S5, 0GPU and 00S1. Full theorem/proof text
was inspected, not only search snippets. No immutable source revision
was pinned or downloaded. The use is a classical specialization, not a
novelty search or a claim of a new obstruction theory. No broader paper
or source cache was changed.

\hypertarget{native:a-topological-obstruction-to-the-proposed-spine-formula}{%
\subsection{4. A topological obstruction to the proposed spine
formula}\label{native:a-topological-obstruction-to-the-proposed-spine-formula}}

Consider the recurrence \(a_{n+1}=a_n\) with integer initial value. Its
global parameter space is \(X=\mathbb Z\). Its local parameter spaces
are \(X_p=\mathbb Z_p\), and every local realization has the same
prime-independent recurrence. Give the local factors their usual
\(p\)-adic topologies and the product its product topology. For these
integral local factors, the restricted product relative to all
\(\mathbb Z_p\) is the ordinary product. Write

\[A=\mathbb Z_2,\qquad B=\prod_{p>2}\mathbb Z_p,\qquad D=\{(n,(n)_{p>2}):n\in\mathbb Z\}\subset A\times B.\]

\textbf{Theorem.} There are no Hausdorff space \(C\) and continuous maps
\(f:A\to C\), \(g:B\to C\) for which the equality locus

\[\{(a,b):f(a)=g(b)\}\]

is exactly \(D\). Thus the canonical local realizations of global
integer constant sequences cannot satisfy the manuscript's fiber-product
formula in this Hausdorff topological interpretation.

\textbf{Proof.} The equality locus is closed: it is the inverse image of
the diagonal of \(C\times C\) under the continuous map \(f\times g\),
and a Hausdorff space has closed diagonal.

The set \(D\) is dense in \(A\times B\). A nonempty basic open set
prescribes congruences at only finitely many primes and finite precision
at each. The moduli are pairwise coprime. The finite Chinese remainder
theorem supplies an integer meeting all those congruences, so every such
open set meets \(D\).

However, \(D\) is not the whole product. The element with 2-adic
coordinate zero and every odd-adic coordinate one is not diagonal. If an
integer had zero 2-adic coordinate, it would be divisible by every power
of two, forcing it to be zero. Its 3-adic coordinate would then be zero,
not one. Consequently \(D\) is a proper dense subset and is not closed.
It cannot be the equality locus above. \(\square\)

This is not asserted to refute the original, underspecified sentence
under every possible reading. It refutes the explicit and natural
interpretation just stated. The assumptions of continuity, Hausdorffness
and identification through the canonical local maps are essential to
this claim.

\hypertarget{native:a-precise-reconstruction-that-does-hold}{%
\subsubsection{A precise reconstruction that does
hold}\label{native:a-precise-reconstruction-that-does-hold}}

For a finite set \(S\) of primes, positive precisions \(L_p\), and a
fixed integer polynomial system \(F\), put

\[N=\prod_{p\in S}p^{L_p},\qquad\operatorname{Sol}_F(R)=\{x\in R^d:F(x)=0\}.\]

Then coefficientwise CRT gives a canonical bijection

\[\operatorname{Sol}_F(\mathbb Z/N\mathbb Z)\cong\prod_{p\in S}\operatorname{Sol}_F(\mathbb Z/p^{L_p}\mathbb Z).\]

The reason is exact: the CRT ring isomorphism preserves polynomial
evaluation, and zero in a finite product is equivalent to zero in every
component. For \(F\) empty, passing to all compatible finite congruences
identifies the profinite completion \(\widehat{\mathbb Z}\) with
\(\prod_p\mathbb Z_p\). The same finite reconstruction is compatible
with coefficient projections and integer recurrence matrices because
these operations commute with reduction and CRT.

Here the polynomial equations or recurrence matrices are fixed common
data. The theorem reconstructs finite or profinite objects. It does not
identify a ``smallest'' spine, impose additional spectral descent, or
recover global \(\mathbb Z\) from arbitrary tuples of local points. A
completed version of 13.3 must say exactly which object is completed and
specify its universal property.

\hypertarget{native:a-finite-compatibility-criterion-with-a-precise-universal-property}{%
\subsubsection{A finite compatibility criterion with a precise universal
property}\label{native:a-finite-compatibility-criterion-with-a-precise-universal-property}}

Fix finite sets \(A,B\) of local packets and an explicitly given
relation \(R\subseteq A\times B\). Say that \(R\) is rectangular when

\[R(a,b),\ R(a',b),\ R(a',b')\quad\Longrightarrow\quad R(a,b')\]

for all \(a,a'\in A\) and \(b,b'\in B\).

\textbf{Theorem.} The relation \(R\) can be represented exactly as

\[R(a,b)\quad\Longleftrightarrow\quad f(a)=g(b)\]

for maps from \(A\) and \(B\) to a common set if and only if \(R\) is
rectangular. In that case a canonical finite choice of common set is the
connected-component set \(C_R\) of the bipartite graph with vertex set
\(A\sqcup B\) and edges given by \(R\).

The component maps \(f_R:A\to C_R\) and \(g_R:B\to C_R\) have the
following universal property. For every set \(Y\) and maps \(h:A\to Y\),
\(k:B\to Y\) satisfying \(h(a)=k(b)\) whenever \(R(a,b)\), there is a
unique map \(u:C_R\to Y\) with \(u\circ f_R=h\) and \(u\circ g_R=k\).

\textbf{Proof.} Any equality representation is rectangular, since the
three assumed equalities imply \(f(a)=g(b)=f(a')=g(b')\).

Conversely, suppose \(R\) is rectangular. An edge always gives equal
component labels. If \(a\in A\) and \(b\in B\) have equal component
labels, there is a finite path between them. Its length is odd. A path
of length one is already an \(R\)-edge. A longer path begins with three
edges \(a-b_1-a_1-b_2\). Rectangularity gives the edge \(a-b_2\).
Replacing these three edges by that edge shortens the path by two.
Induction therefore yields \(R(a,b)\). This proves the exact equality
representation. Finiteness follows because the component set is a
quotient of the finite set \(A\sqcup B\).

For the universal property, combine \(h\) and \(k\) into a function on
\(A\sqcup B\). It is constant across every edge by hypothesis, hence
along every finite path. It therefore descends to a well-defined
function on components. The two required factorization identities hold
by construction. They determine the value on every component, proving
uniqueness. \(\square\)

For example, the three pairs \((a_0,b_0),(a_1,b_0),(a_1,b_1)\) without
\((a_0,b_1)\) fail rectangularity. Every equality-label representation
agreeing on those three pairs also admits the fourth pair, so it cannot
represent that relation exactly.

This is a precise finite gluing obligation: future arithmetic
applications must define \(A\), \(B\) and \(R\) and prove
rectangularity. The universal property holds for the component quotient
even when rectangularity fails, but then the equality locus strictly
enlarges \(R\). The theorem establishes neither a minimum number of
labels nor a smallest spine in the manuscript's unspecified category. It
carries no automatic topology, spectral information or Frobenius
compatibility. \textbf{Evidence: prose proof in this note, not a Lean
formalization.}
